% Part: proof-theory
% Chapter: cut-elimination
% Section: introduction

\documentclass[../../../include/open-logic-section]{subfiles}

\begin{document}

\olfileid[hi]{pt}{cut}{int}

\olsection{परिचय}

\Cut{} नियम यह है
\[
\Axiom$\Gamma \fCenter \Delta, !A$
\Axiom$!A, \Pi \fCenter \Lambda$
\RightLabel{\Cut}
\BinaryInf$\Gamma, \Pi \fCenter \Delta, \Lambda$
\DisplayProof
\]
(!!{formula}~$!A$ को \emph{कट सूत्र} कहते हैं।)

जेंट्सेन के मूल अनुक्रम कलन~\Log{LK} में \Cut{} सम्मिलित था। जेंट्सेन का
\emph{हाउप्टज़ाट्स} (``मुख्य प्रमेय'') कहता है कि~\Log{LK} में !!{provable}
प्रत्येक अनुक्रम \Cut{} नियम का उपयोग किए बिना भी !!{provable} है; अर्थात्
\Log{LK}-प्रमाणों से \Cut{} का ``विलोपन'' किया जा सकता है। हमारी प्रणालियों
\Log{G1c} और \Log{G3c} में \Cut नहीं है, किंतु उनके लिए भी प्रश्न उठाया जा
सकता है: क्या \Cut{} तथा \Log{G1c} (या \Log{G3c}) के अन्य नियमों का उपयोग
करने वाले !!{proof}s से \Cut{} हटाया जा सकता है? हमारी स्थापित शब्दावली में
समतुल्य प्रश्न है: क्या \Cut{}, \Log{G1c} (या \Log{G3c}) का ग्राह्य नियम है?

कट-विलोपन प्रमेय शायद प्रमाण सिद्धांत का सबसे महत्त्वपूर्ण और प्रसिद्ध
परिणाम है। सर्वप्रथम यह स्थापित करता है कि \Log{G1c} और \Log{G3c} को ऐसे
नियम की आवश्यकता नहीं है जो पहली दृष्टि में स्पष्टतः आवश्यक लग सकता है। जब
आप वास्तविक प्रमाणों को औपचारिक रूप देने का प्रयास करेंगे, तो देखेंगे कि
सहायक प्रमेयों के उपयोग से निपटने का कोई ढंग चाहिए। गणितज्ञ परिणाम सिद्ध करते
हैं और फिर दूसरे परिणामों के प्रमाणों में उनका आह्वान करते हैं। अतः आप पहले
परिणाम $L$ (सहायक प्रमेय) के प्रमाण को अनुक्रम $\Sequent L$ के !!{proof} के
रूप में औपचारिक कर सकते हैं। फिर दूसरे परिणाम~$R$, जिसके प्रमाण में~$L$ का
उपयोग है, को अनुक्रम $L \Sequent R$ के !!{proof} के रूप में औपचारिक करते हैं।
आप कैसे जानेंगे कि~$R$ का प्रमाण, अर्थात् केवल $\Sequent R$ का !!a{proof},
है? दोनों !!{proof}s को मिलाने का कोई ढंग चाहिए, और \Cut{} नियम ऐसा करने देता
है।

हमने बताया कि \Log{G1c} और \Log{G3c} पूर्ण हैं, इसलिए वे हर उस अनुक्रम को
सिद्ध करते हैं जिसे सिद्ध हो सकने की अपेक्षा होगी—सभी वैध अनुक्रमों को। इसे
सीधे सिद्ध किया जा सकता है। किंतु जेंट्सेन के समय केवल अभिगृहीतीय व्युत्पत्ति
प्रणालियों की पूर्णता ज्ञात थी। \Log{LK} की पूर्णता दिखाने के लिए जेंट्सेन ने
अभिगृहीतीय प्रणाली की !!{derivation}s का \Log{LK}-!!{proof}s में अनुवाद दिया।
इस अनुवाद में \Cut{} का अनिवार्य उपयोग था। कट-विलोपन केवल चिरसम्मत तर्क का
महत्त्वपूर्ण परिणाम नहीं है: अनेक अन्य तर्कों के भी अनुक्रम कलन हैं। कुछ
तर्कों के अनुक्रम कलनों की पूर्णता का सीधा प्रमाण कठिन या असंभव है, इसलिए अन्य
प्रणालियों से अनुवाद (\Cut का उपयोग करके) ही पूर्णता दिखाने का एकमात्र मार्ग
है। तब यह स्थापित करने के लिए प्रमाण-सैद्धांतिक ढंग से कट-विलोपन सिद्ध करना
आवश्यक है कि \Cut{} अनावश्यक है और \Cut{}-विहीन प्रणाली पूर्ण है।

कट-विलोपन इसलिए भी महत्त्वपूर्ण है कि उसके अनुप्रयोग से अन्य स्वतंत्र रूप से
रोचक परिणाम मिलते हैं। हम दो उदाहरणों—क्रेग अंतर्वेशन प्रमेय और हेरब्रांड
प्रमेय—पर विचार करेंगे।

कट-विलोपन सिद्ध करने में प्रयुक्त विचार सामान्यतः यह दिखाते हैं कि \Cut{} का
उपयोग करने वाले !!a{proof} को उसके बिना वाले प्रमाण में कैसे बदला जाए। ये
रूपांतरण प्रायः ``स्थानीय'' होते हैं; अर्थात् हम किसी !!a{proof} का एक भाग
(सामान्यतः \Cut{} अनुमान पर समाप्त होने वाला उप-!!{proof}) लेकर उसे उसी
अनुक्रम के दूसरे उप-!!{proof} से बदलते हैं, जिसमें \Cut{} अनुमान हटा दिया गया
हो या अन्य अनुमानों से बदल दिया गया हो। ऐसे तर्क \Cut{} वाले !!{proof}s को
उसके बिना वाले प्रमाणों में बदलने की चरण-दर-चरण कलन-विधियाँ देते हैं।
अनुप्रयोगों में कट-विलोपन प्रक्रियाओं के उपयोग पर इससे अधिक सूचना मिलती है।
उदाहरणतः अंतर्वेशन प्रमेय के निदर्श-सैद्धांतिक प्रमाण इस रूप के हैं: ``यदि $!A
\lif !B$ वैध है तो कोई अंतर्वेशक विद्यमान है'' (अभी इस बात को छोड़ दें कि
अंतर्वेशक क्या है)। कट-विलोपन वाला प्रमाण-सैद्धांतिक तर्क ऐसी कलन-विधि देता है
जो $!A \lif !B$ का !!a{proof} लेकर अंतर्वेशक लौटाती है। निदर्श-सैद्धांतिक तर्क
के विपरीत प्रमाण-सैद्धांतिक तर्क \emph{रचनात्मक} है।

कट-विलोपन प्रमेय सिद्ध करने के दो लोकप्रिय ढंग हैं। एक (जेंट्सेन तक जाने
वाला) दिखाता है कि प्रमाण के सबसे ऊपर के कट हटाए जा सकते हैं, और फिर ऐसे कटों
को तब तक हटाता है जब तक कोई शेष न रहे। दूसरा (मूलतः श्यूटे और टेट का) किसी
!!{proof} के सबसे जटिल कटों पर केंद्रित होता है और क्रमशः उन कटों को हटाता है
जो इस अर्थ में महत्तम हैं कि किसी अन्य कट के कट !!{formula} की गहराई अधिक
नहीं है। दोनों के अपने लाभ और कमियाँ हैं। कुछ प्रणालियों में एक उपागम काम
करता है जबकि दूसरे का कार्यान्वयन कठिन या असंभव भी होता है। इसलिए हम दोनों
उपागमों का वर्णन करते हैं। हम~\Log{G3c} के लिए कट-विलोपन सिद्ध करेंगे। वास्तव
में हम \Cut{} का नहीं, बल्कि \emph{प्रसंग-साझा कट}~\CutCS का विलोपन दिखाएँगे:
\[
\Axiom$\Gamma \fCenter \Delta, !A$
\Axiom$!A, \Gamma \fCenter \Delta$
\RightLabel{\Cut}
\BinaryInf$\Gamma \fCenter \Delta$
\DisplayProof
\]
यदि कोई अनुक्रम कलन दुर्बलीकरण और संकुचन स्वीकार करता है (या तो \Log{G1c}
जैसे वास्तविक नियमों के रूप में या \Log{G3c} जैसे ग्राह्य नियमों के रूप में),
तो \Cut{}, \CutCS का अनुकरण कर सकता है और इसका विपरीत भी। हम \Cut{} के बजाय
\CutCS{} चुनते हैं, क्योंकि पहली कट-विलोपन युक्ति का~\CutCS के लिए वर्णन सरल
है और दूसरी युक्ति केवल~\CutCS के लिए काम करती है।

\end{document}
